The v5 atlas contains 20,000 labeled datasets split into 10,000 discovery and 10,000 validation rows. They are selected from a 50,000-row candidate pool and span 52 base generators grouped into 16 broad families. The families include chaotic maps and flows, input-driven memory systems, oscillatory and quasiperiodic systems, long-range processes, colored noise, drifting trends, regime switching, seasonal/calendar structure, volatility-heteroskedastic processes, delay dynamics, intermittent sparse processes, nonlinear autoregressions, heavy-tail jumps, and multiscale composites.

The 50,000-row candidate pool is used to make the final atlas broad without turning the paper into an unstructured random sweep. Candidate series are first generated across the full taxonomy with varied seeds and parameters. The final discovery and validation rows are then selected to cover both broad family balance and frontier enrichment near the regions where QRC and ESN are close. This matters scientifically: a purely average benchmark would mostly estimate that ESN is the stronger default, whereas a regime atlas must also contain enough near-boundary and QRC-favorable examples to learn where the boundary lies.

\begin{table}[!t]
\centering
\caption{Frozen v5 protocol and held-out validation summary.}
\label{tab:v5-summary}
\footnotesize
\setlength{\tabcolsep}{3.0pt}
\renewcommand{\arraystretch}{1.03}
\begin{tabularx}{\columnwidth}{@{}lX@{}}
\toprule
Item & Value \\
\midrule
Discovery / validation rows & 10,000 / 10,000 \\
Families / base generators & 16 / 52 \\
Measured descriptors & 30 Tier-A properties \\
QRC variants & QRC-M, QRC-E, QRC-D \\
Feature dimension & 63 for QRC and ESN \\
Best-QRC win rate & 32.4\% on validation \\
Best-QRC useful rate & 12.5\% on validation \\
Mean best-QRC advantage & \(-0.058\) on validation \\
\bottomrule
\end{tabularx}
\end{table}

The target is one-step or short-horizon forecasting with train/validation/test splits derived from each generated series. All preprocessing is fit on training data only. The v5 evaluation records NRMSE and NMAE for linear ridge, gradient boosting, NVAR, the frozen ESN, and all three QRC variants. The primary label is the NRMSE advantage of the best QRC variant over the ESN. We also retain variant-specific advantages so that the atlas can distinguish whether QRC-M, QRC-E, or QRC-D is responsible for a QRC-favorable row.

\begin{table}[!t]
\centering
\caption{Descriptor blocks used for the 30-feature regime map.}
\label{tab:feature-blocks}
\scriptsize
\setlength{\tabcolsep}{3.0pt}
\renewcommand{\arraystretch}{1.05}
\begin{tabularx}{\columnwidth}{@{}lX@{}}
\toprule
Block & Examples \\
\midrule
Memory and persistence & Autocorrelation time, memory capacity, DFA exponent, Hurst estimate \\
Predictability & Linear and GBM NRMSE, nonlinear-gain contrast, predictability gap \\
Spectrum & Spectral entropy, flatness, dominant frequency, spectral slope and centroid \\
Drift and stationarity & ADF/KPSS signals, differencing order, trend strength, changepoints \\
Noise and volatility & SNR, volatility autocorrelation, ARCH statistic \\
Nonlinear structure & Lyapunov estimate, 0--1 statistic, entropy, recurrence, Lempel--Ziv complexity \\
\bottomrule
\end{tabularx}
\end{table}

NRMSE is the primary metric because it is scale-free and standard for forecasting comparisons. NMAE is recorded as a robustness metric because it is less sensitive to rare large errors. The main paper uses NRMSE labels; the artifact tables preserve NMAE columns for follow-up robustness checks and for readers who prefer absolute-error normalization.

The paper-facing analysis package writes all tables and figures from the completed v5 outputs. The main validation outputs are the complete validation table, the split/family/variant summaries, the discovery-trained validation predictions, the feature-set robustness table, and the rule-pocket table. All reported numbers in this paper are taken from these generated artifacts.
